\section{Vector Indexes and Filtered Search}
\label{sec-vectors}

\sys keeps one unit-length float vector $\mathbf{v}(t)$ per term in memory,
loaded from a checksummed file whose records name terms by spelling, so a
file survives a dictionary that numbers terms differently.
Embeddings come from outside; \sys stores and searches them.
Four index families are available, each a different trade between
exactness, memory, and build time.

\stitle{Flat}.
Score every vector: $O(Nd)$ and exact.
With 8-bit scalar quantization (one scale $\delta=\max_i|v_i|/127$ per vector, for coordinates $v_i$) each
coordinate is off by at most $\delta/2$, so for a unit query a score is off by at
most $(\delta/2)\sqrt d$, in a quarter of the memory.

\stitle{IVF}.
Spherical $k$-means with $k$-means++ seeding~\cite{arthur2007} gives $\nu$
centroids and each vector goes to its nearest.
A query scores the centroids and reads the $\rho$ nearest lists.
The cost is $\nu d+\rho (N/\nu) d$, minimal at $\nu=\sqrt{\rho N}$.
The design choice is where the lists live: \emph{as key ranges of the B+ tree
that holds the triples}.
The pair (list, row) is a key $(\mathrm{list},\mathrm{row},0)$, so a list is a
prefix scan, an insert is a tree insert, and the lists page, cache, and latch
like every other index of the engine.
Under a filter the search keeps probing the nearest remaining lists until it
holds $k$ allowed rows, so a filter that excludes everything near the query
costs more instead of returning nothing.

\stitle{Product quantization (PQ)}.
Cut a vector into $m$ pieces and replace each by the index of its nearest of
256 learned centers~\cite{jegou2011}: $m$ bytes per vector.
A query builds $m$ tables of 256 dot products and scores a code as the sum of
$m$ lookups, which equals the dot product with the decoded vector, so all the
error is on the stored side.

\stitle{HNSW}.
The layered proximity graph of Malkov and Yashunin~\cite{malkov2020} with the
paper's neighbor-selection heuristic: a candidate is kept only if it is
nearer the new node than to every neighbor already kept.
Layer 0 allows $2M$ links, and the search beam is $\ef$.
Search takes an optional predicate and still crosses nodes it may not return.

\subsection{Filtering a Similarity Query by a Class}
\label{sec-hybrid}

A filtered similarity query $\mathrm{near}(q,c,k)$ asks for the $k$ terms of
$\desc(c)$ nearest to $q$.
Let $f(c)=|\desc(c)|/N$.
We compare four ways to answer it, all returning terms of $\desc(c)$ only.
\begin{description}
\item[Pre-filter.] Resolve $c$ to its members first (one range scan per
  interval of $I(c)$), then score only those, exactly. About $f(c)\,N\,d$ work.
\item[Post-filter.] Ask the index for $k'>k$ terms, keep those below $c$ (a
  membership test is a binary search over $I(c)$), and multiply $k'$ by four
  until $k$ remain or the index is exhausted.
  If the filter were independent of the similarity, about $k/f(c)$
  candidates would be needed.
\item[Masked search.] Give the index a predicate, ``the row's number lies in
  an interval of $I(c)$'', so a list scan or a graph walk skips rows outside
  $c$.
  Nothing is resolved in advance: each row the search touches costs one
  binary search.
  HNSW still crosses rows it may not return, to keep the graph connected, in
  the spirit of ACORN~\cite{patel2024}; IVF keeps probing lists until it has
  $k$ allowed rows.
\item[Auto.] Pre-filter when $|\desc(c)|$ is at most 4,096, otherwise masked
  HNSW.
\end{description}

\stitle{The filter is correlated with the vectors}.
In an ontology the names of a class's descendants are close to the class and
to each other, so the neighbors of a term tend to lie under the term's own
ancestors.
On the Gene Ontology, of the unfiltered top 100 neighbors of a query from
\emph{inside} $c$, 4\% to 72\% lie under $c$, rising with the size of $c$; for
a query from \emph{outside} $c$ the share is 0\% to 12\%.
The independence estimate $k/f(c)$ is therefore wrong in both directions:
post-filtering is cheap for an inside query and does not terminate usefully
for an outside one.

\stitle{Do not resolve the filter}.
A masked search that first resolved $c$ by interval scans and built a bitmap
per query spent most of its time resolving: for a class of over 5,000 members
that was 560 of 669 $\mu$s.
Testing the number of each row the search touches, one binary search over a
handful of intervals, brought the same search to 105 $\mu$s.
This is where interval labels earn their keep a second time: they make
membership an $O(\log|I(c)|)$ test with no scan.
